% ============================================================
%  GUÍA 5: QUÍMICA — MATERIA, MEZCLAS Y SEPARACIONES
%  Curso Introductorio — 3er Año
%  Autor: Ing. Gabriel Astudillo
%  Sesiones cubiertas: 13, 14, 17, 18
% ============================================================

\documentclass[12pt, a4paper]{article}

% ── Paquetes ─────────────────────────────────────────────────

\usepackage{mathwizards-guia}
\mwguia{Guía 5 — Química: Materia, Mezclas y Separaciones}{Ing. Gabriel Astudillo $\cdot$ Curso 3er Año}

\begin{document}
% ============================================================

% ── Portada / Título ─────────────────────────────────────────
\mwportada{Guía de Estudio 5}{Química: Materia, Mezclas y Separaciones}{Método Científico, Laboratorio, Propiedades de la Materia y Técnicas de Separación}

\vspace{4pt}
\begin{cuadroinfo}
\small
\textbf{Presentación:} La Química es la ciencia que estudia de qué están hechas las cosas y cómo se transforman. En esta guía descubrirás cómo trabajan los científicos (el método científico), conocerás las herramientas del laboratorio, aprenderás a clasificar la materia que nos rodea y dominarás las técnicas para separar mezclas. Todo esto con un ojo analítico: no se trata de memorizar, sino de entender \emph{por qué} cada método funciona.
\\[4pt]
\textbf{Nivel de Dificultad:}\quad
\facil{} Básico / Directo \quad
\medio{} Intermedio \quad
\dificil{} Desafiante
\end{cuadroinfo}

\vspace{8pt}

% ============================================================
\section{El Método Científico y el Laboratorio}
% ============================================================

\subsection{El método científico}

El \textbf{método científico} es el procedimiento sistemático que usan los científicos para investigar fenómenos naturales, formular hipótesis y validar sus conclusiones. No es una receta rígida, sino una guía flexible basada en la evidencia.

\begin{cuadroteoria}
\textbf{Etapas del método científico:}
\begin{enumerate}[leftmargin=*, label=\arabic*.]
  \item \textbf{Observación:} Percibir un fenómeno o hecho que despierta curiosidad.
  \item \textbf{Planteamiento del problema:} Formular una pregunta clara y específica.
  \item \textbf{Formulación de hipótesis:} Proponer una explicación provisional y comprobable.
  \item \textbf{Experimentación:} Diseñar y ejecutar experimentos controlados para poner a prueba la hipótesis.
  \item \textbf{Análisis de resultados:} Organizar y examinar los datos obtenidos.
  \item \textbf{Conclusión:} Determinar si los datos apoyan o refutan la hipótesis.
  \item \textbf{Comunicación:} Publicar los resultados para que otros puedan revisarlos y reproducirlos.
\end{enumerate}
\end{cuadroteoria}

\begin{cuadroadvertencia}
\textbf{Breve historia:} La Química nació de la \emph{alquimia}, una práctica medieval que buscaba transformar metales comunes en oro. Aunque los alquimistas nunca lo lograron, sus experimentos sentaron las bases de la Química moderna. Antoine Lavoisier (siglo XVIII) es considerado el «padre de la Química» por establecer la Ley de Conservación de la Masa y desterrar la teoría del flogisto.
\end{cuadroadvertencia}

\subsection{El instrumental de laboratorio}

En un laboratorio de Química, cada instrumento tiene un propósito específico. Se clasifican según su función:

\begin{cuadrosolucion}
\textbf{Clasificación del instrumental:}
\begin{itemize}[leftmargin=*]
  \item \textbf{De sostén y soporte:} Permiten sujetar otros instrumentos. Ejemplos: soporte universal, pinzas, aros metálicos, trípode, gradilla.
  \item \textbf{De uso específico:} Diseñados para funciones particulares. Ejemplos: mechero de Bunsen (calentar), termómetro (medir temperatura), balanza (medir masa), embudo (filtrar/trasvasar).
  \item \textbf{Volumétricos:} Miden volúmenes de líquidos con precisión. Ejemplos: probeta graduada, bureta, pipeta, matraz aforado.
  \item \textbf{Recipientes:} Contienen sustancias. Ejemplos: vaso de precipitados (beaker), matraz Erlenmeyer, tubo de ensayo, vidrio de reloj, cápsula de porcelana.
\end{itemize}
\end{cuadrosolucion}



\subsection{Normas de seguridad en el laboratorio}

\begin{cuadroextra}
\textbf{Reglas fundamentales:}
\begin{enumerate}[leftmargin=*, label=\arabic*.]
  \item Usa siempre bata de laboratorio, lentes de seguridad y guantes cuando sea necesario.
  \item Nunca pruebes, huelas directamente ni toques sustancias químicas sin indicación del profesor.
  \item Para percibir olores, usa la técnica del \textbf{abanico}: agita suavemente con la mano los vapores hacia tu nariz.
  \item Nunca calientes un recipiente cerrado.
  \item Lee las etiquetas de los reactivos \textbf{antes} de usarlos.
  \item Conoce la ubicación del extintor, la ducha de emergencia y el botiquín.
  \item Al calentar un tubo de ensayo, apunta su boca \textbf{lejos} de ti y de los demás.
\end{enumerate}
\end{cuadroextra}

\newpage

% ============================================================
\subsection{Ejercicios: Método Científico y Laboratorio}
% ============================================================

\begin{enumerate}[leftmargin=*]
  \item \facil{} Ordena las siguientes etapas del método científico: Experimentación, Observación, Conclusión, Hipótesis, Planteamiento del problema.

  \item \facil{} ¿Cuál es la diferencia entre una observación y una hipótesis? Da un ejemplo de cada una.

  \item \facil{} Clasifica cada instrumento en su categoría (sostén, uso específico, volumétrico o recipiente):
  \begin{enumerate}[label=\alph*)]
    \item Probeta graduada
    \item Soporte universal
    \item Tubo de ensayo
    \item Mechero de Bunsen
  \end{enumerate}

  \item \facil{} Nombra dos normas de seguridad que deben seguirse al trabajar con sustancias químicas en el laboratorio.

  \item \medio{} Un estudiante observa que las plantas en la ventana crecen más que las que están en el interior. Formula una hipótesis y diseña un experimento sencillo para comprobarla. Identifica la variable independiente, la variable dependiente y las variables controladas.

  \item \medio{} ¿Qué instrumento utilizarías en cada caso?
  \begin{enumerate}[label=\alph*)]
    \item Medir exactamente $25\;\text{mL}$ de una solución.
    \item Calentar una sustancia líquida.
    \item Pesar $50\;\text{g}$ de sal.
    \item Filtrar una mezcla de arena y agua.
  \end{enumerate}

  \item \medio{} Explica por qué es incorrecto husmear directamente un frasco de reactivo y describe la técnica correcta.

  \item \medio{} ¿Por qué Lavoisier es considerado el padre de la Química moderna? ¿Qué ley fundamental estableció?

  \item \dificil{} Un grupo de estudiantes quiere determinar si la temperatura del agua afecta la velocidad con que se disuelve el azúcar. Diseña el experimento completo indicando: hipótesis, materiales, procedimiento paso a paso, variables y cómo registrarían los datos.

  \item \dificil{} ¿Es posible que una hipótesis resulte falsa y aun así el experimento sea exitoso desde el punto de vista científico? Justifica tu respuesta.

  \item \dificil{} Un científico afirma: «He probado mi hipótesis y es definitivamente verdadera para siempre.» ¿Esta afirmación es compatible con el método científico? Explica el concepto de \emph{falsabilidad}.

  \item \dificil{} Se tiene un líquido desconocido en un laboratorio. Diseña un plan de análisis (usando solo los instrumentos y técnicas vistas en esta guía) para determinar: su volumen, su masa, su densidad y si es una sustancia pura o una mezcla. Justifica la elección de cada instrumento.
\end{enumerate}

\newpage

% ============================================================
\section{La Materia y sus Propiedades}
% ============================================================

\subsection{Definición de materia}

\begin{cuadroteoria}
\textbf{Materia} es todo aquello que tiene \textbf{masa} y ocupa un \textbf{lugar en el espacio} (volumen). Desde una estrella hasta un átomo, todo lo que puedes tocar (y mucho de lo que no puedes) es materia.
\end{cuadroteoria}

Un \textbf{material} es una porción específica de materia que podemos estudiar. Ejemplos: un trozo de madera, un vaso de agua, el aire de una habitación.

\subsection{Propiedades de la materia}

Las propiedades de la materia se dividen en dos grandes grupos:

\begin{cuadrosolucion}
\begin{itemize}[leftmargin=*]
  \item \textbf{Propiedades extensivas} (no características): \textbf{Dependen de la cantidad} de materia. Si tienes más material, estas propiedades cambian.
  \begin{itemize}
    \item Masa, volumen, peso, longitud, número de moles.
  \end{itemize}
  
  \item \textbf{Propiedades intensivas} (características): \textbf{No dependen de la cantidad} de materia. Sirven para \textbf{identificar} una sustancia.
  \begin{itemize}
    \item Densidad, temperatura de ebullición, temperatura de fusión, color, olor, solubilidad, conductividad eléctrica.
  \end{itemize}
\end{itemize}
\end{cuadrosolucion}

\begin{cuadroadvertencia}
\textbf{Ejemplo clave:} Un litro de agua y un vaso de agua tienen diferente masa y volumen (propiedades extensivas), pero ambos tienen la misma densidad ($1\;\text{g/cm}^3$), el mismo punto de ebullición ($100\;°\text{C}$) y el mismo punto de fusión ($0\;°\text{C}$) — propiedades intensivas que nos dicen que ambos son agua.
\end{cuadroadvertencia}

\begin{figure}[h!]
  \centering
  \begin{tikzpicture}[x=1cm, y=1cm]
    % Container 1: Whole sample
    \begin{scope}[shift={(-3.5, 0)}]
      % Liquid
      \fill[mwprimario!12] (-1.2, -1.0) rectangle (1.2, 0.8);
      % Beaker lines
      \draw[very thick, gray] (-1.2, 1.2) -- (-1.2, -1.0) -- (1.2, -1.0) -- (1.2, 1.2);
      
      % Labels
      \node[above, font=\bfseries\small, text=mwprimario] at (0, 1.3) {Muestra Completa};
      \node[align=left, font=\footnotesize] at (0, -2.5) {
        \textbf{Propiedades Extensivas:} \\
        Masa ($m$) = $100\text{ g}$ \\
        Volumen ($V$) = $100\text{ mL}$ \\[4pt]
        \textbf{Propiedades Intensivas:} \\
        Densidad ($\rho$) = $1\text{ g/cm}^3$ \\
        Temperatura ($T$) = $20\,^\circ\text{C}$
      };
    \end{scope}

    % Division arrow
    \draw[ultra thick, -Latex, gray] (-1.6, 0) -- node[above, font=\scriptsize\bfseries, text=gray] {Dividir muestra} (1.6, 0);

    % Container 2 & 3: Half samples
    \begin{scope}[shift={(3.5, 0)}]
      % Sub-beaker 1
      \begin{scope}[shift={(-1.2, 0)}]
        \fill[mwprimario!12] (-0.7, -1.0) rectangle (0.7, -0.1);
        \draw[thick, gray] (-0.7, 0.5) -- (-0.7, -1.0) -- (0.7, -1.0) -- (0.7, 0.5);
        \node[above, font=\bfseries\scriptsize, text=mwprimario] at (0, 0.6) {Muestra A};
        \node[align=left, font=\scriptsize] at (0, -2.3) {
          $m = 50\text{ g}$\\
          $V = 50\text{ mL}$\\
          $\rho = 1\text{ g/cm}^3$\\
          $T = 20\,^\circ\text{C}$
        };
      \end{scope}

      % Sub-beaker 2
      \begin{scope}[shift={(1.2, 0)}]
        \fill[mwprimario!12] (-0.7, -1.0) rectangle (0.7, -0.1);
        \draw[thick, gray] (-0.7, 0.5) -- (-0.7, -1.0) -- (0.7, -1.0) -- (0.7, 0.5);
        \node[above, font=\bfseries\scriptsize, text=mwprimario] at (0, 0.6) {Muestra B};
        \node[align=left, font=\scriptsize] at (0, -2.3) {
          $m = 50\text{ g}$\\
          $V = 50\text{ mL}$\\
          $\rho = 1\text{ g/cm}^3$\\
          $T = 20\,^\circ\text{C}$
        };
      \end{scope}
    \end{scope}
  \end{tikzpicture}
  \caption{Propiedades extensivas (dependen de la cantidad de materia) vs. intensivas (no dependen de ella).}
  \label{fig:propiedades}
\end{figure}

\subsection{Masa vs. Peso}

Estos dos conceptos se confunden frecuentemente, pero son muy diferentes:

\begin{cuadroteoria}
\begin{itemize}[leftmargin=*]
  \item \textbf{Masa ($m$):} Cantidad de materia que posee un cuerpo. Es \textbf{constante} sin importar dónde estés (en la Tierra, en la Luna, en el espacio). Se mide en \textbf{kilogramos} (kg). Es una magnitud \textbf{escalar}.
  
  \item \textbf{Peso ($W$):} Es la \textbf{fuerza} con que la gravedad atrae al cuerpo. \textbf{Varía} según la gravedad del lugar. Se calcula como $W = m \cdot g$. Se mide en \textbf{newtons} (N). Es una magnitud \textbf{vectorial}.
\end{itemize}
\end{cuadroteoria}

Un astronauta de $70\;\text{kg}$ de masa pesa $686\;\text{N}$ en la Tierra ($g = 9{,}8\;\text{m/s}^2$) pero solo $113{,}4\;\text{N}$ en la Luna ($g = 1{,}62\;\text{m/s}^2$). ¡Su masa sigue siendo la misma!

%\newpage

% ============================================================
\subsection{Ejercicios: Materia y Propiedades}
% ============================================================

\begin{enumerate}[leftmargin=*, resume]
  \item \facil{} Clasifica cada propiedad como extensiva (E) o intensiva (I):
  \begin{multicols}{2}
  \begin{enumerate}[label=\alph*)]
    \item Masa
    \item Densidad
    \item Volumen
    \item Punto de ebullición
    \item Color
    \item Peso
  \end{enumerate}
  \end{multicols}

  \item \facil{} ¿Cuál es la diferencia fundamental entre masa y peso?

  \item \facil{} Un objeto tiene una masa de $5\;\text{kg}$. Calcula su peso en la Tierra ($g = 9{,}8\;\text{m/s}^2$) y en la Luna ($g = 1{,}62\;\text{m/s}^2$).

  \item \facil{} ¿Por qué la densidad es una propiedad intensiva? Explica con un ejemplo.

  \item \medio{} Un estudiante tiene dos trozos de un metal desconocido. El primero tiene una masa de $100\;\text{g}$ y un volumen de $37\;\text{cm}^3$. El segundo tiene una masa de $250\;\text{g}$ y un volumen de $92{,}6\;\text{cm}^3$. ¿Son del mismo material? Justifica calculando.

  \item \medio{} Verdadero o falso. Justifica cada respuesta:
  \begin{enumerate}[label=\alph*)]
    \item «La masa de un cuerpo cambia si lo llevamos a la Luna.»
    \item «El punto de ebullición del agua siempre es $100\;°\text{C}$.»
    \item «El volumen es una propiedad que sirve para identificar una sustancia.»
    \item «El peso de un objeto se mide en kilogramos.»
  \end{enumerate}

  \item \medio{} Si cortas un bloque de hierro por la mitad, ¿qué propiedades cambian y cuáles se mantienen igual? Menciona al menos dos de cada tipo.

  \item \medio{} ¿Es el aire materia? Justifica tu respuesta usando la definición de materia.

  \item \dificil{} Un cuerpo pesa $147\;\text{N}$ en la Tierra. ¿Cuánto pesa en un planeta donde la gravedad es $3{,}5\;\text{m/s}^2$? Primero calcula su masa.

  \item \dificil{} Explica por qué en la vida cotidiana decimos «peso 70 kilos» aunque técnicamente es incorrecto. ¿Qué deberíamos decir realmente?

  \item \dificil{} Se descubre un nuevo material en un laboratorio. Se toman tres muestras de diferentes tamaños y se miden las siguientes propiedades:
  \begin{center}
  \small
  \begin{tabular}{lccccc}
  \toprule
  & Masa (g) & Volumen (cm$^3$) & Densidad & Color & Punto de fusión \\
  \midrule
  Muestra 1 & 50 & 18{,}5 & ? & Plateado & $660\;°$C \\
  Muestra 2 & 120 & 44{,}4 & ? & Plateado & $660\;°$C \\
  Muestra 3 & 200 & 74{,}1 & ? & Plateado & $660\;°$C \\
  \bottomrule
  \end{tabular}
  \end{center}
  Calcula la densidad de cada muestra. ¿Confirman que se trata del mismo material? ¿De qué material crees que se trata?

  \item \dificil{} Un astronauta pesa $784\;\text{N}$ en la Tierra. Si viaja a Marte ($g_{\text{Marte}} = 3{,}72\;\text{m/s}^2$) y luego a Júpiter ($g_{\text{Júpiter}} = 24{,}79\;\text{m/s}^2$), calcula su peso en cada planeta. ¿En cuál pesaría más que en la Tierra?
\end{enumerate}

%\newpage

% ============================================================
\section{Taxonomía de la Materia: Sustancias y Mezclas}
% ============================================================

\subsection{Clasificación general de la materia}

Toda la materia que nos rodea puede clasificarse en un esquema jerárquico según su composición:

\begin{cuadroteoria}
\textbf{Clasificación de la materia:}
\begin{itemize}[leftmargin=*]
  \item \textbf{Sustancias puras:} Tienen composición definida y constante.
  \begin{itemize}
    \item \textbf{Elementos:} No se pueden descomponer en sustancias más simples por métodos químicos. Ej: oxígeno (O$_2$), hierro (Fe), oro (Au).
    \item \textbf{Compuestos:} Formados por dos o más elementos combinados en proporciones fijas. Se pueden descomponer químicamente. Ej: agua (H$_2$O), sal (NaCl), dióxido de carbono (CO$_2$).
  \end{itemize}
  
  \item \textbf{Mezclas:} Combinación de dos o más sustancias \textbf{sin} reacción química. Cada componente mantiene sus propiedades.
  \begin{itemize}
    \item \textbf{Mezclas homogéneas} (soluciones): No se distinguen sus componentes a simple vista. Tienen \textbf{una sola fase}. Ej: agua con sal (salmuera), aire, bronce (aleación Cu-Sn).
    \item \textbf{Mezclas heterogéneas:} Se pueden distinguir sus componentes. Tienen \textbf{dos o más fases}. Ej: ensalada, arena con piedras, agua con aceite.
  \end{itemize}
\end{itemize}
\end{cuadroteoria}

\begin{figure}[h!]
  \centering
  \begin{tikzpicture}[
      level 1/.style={sibling distance=7.5cm, level distance=1.5cm},
      level 2/.style={sibling distance=3.5cm, level distance=1.8cm},
      root/.style={draw=mwprimario, fill=mwprimario!12, thick, rounded corners, inner sep=6pt, font=\bfseries\small, text=mwprimario},
      node1/.style={draw=mwmagenta, fill=mwmagenta!12, thick, rounded corners, inner sep=5pt, font=\bfseries\scriptsize, text=mwmagenta, align=center},
      node2/.style={draw=mwnaranja, fill=mwnaranja!12, thick, rounded corners, inner sep=4pt, font=\scriptsize, text=mwnaranja, align=center}
    ]
    % Root
    \node[root] {MATERIA}
      % Sustancias Puras
      child { node[node1] {SUSTANCIAS PURAS\\(Composición fija)}
        child { node[node2] {\textbf{Elementos}\\Ej: Hierro ($\text{Fe}$)\\Oxígeno ($\text{O}_2$)} }
        child { node[node2] {\textbf{Compuestos}\\Ej: Agua ($\text{H}_2\text{O}$)\\Sal ($\text{NaCl}$)} }
      }
      % Mezclas
      child { node[node1] {MEZCLAS\\(Composición variable)}
        child { node[node2] {\textbf{Homogéneas}\\(Una sola fase)\\Ej: Salmuera, Aire} }
        child { node[node2] {\textbf{Heterogéneas}\\(Dos o más fases)\\Ej: Agua y Aceite} }
      };
  \end{tikzpicture}
  \caption{Clasificación jerárquica de la materia.}
  \label{fig:clasificacion}
\end{figure}

\subsection{Clasificación granulométrica de las mezclas}

Las mezclas heterogéneas se pueden subclasificar según el \textbf{tamaño de las partículas} dispersas:

\begin{cuadrosolucion}
\begin{itemize}[leftmargin=*]
  \item \textbf{Dispersiones groseras:} Partículas grandes, visibles a simple vista. Se sedimentan con el tiempo. Ej: agua con arena, jugo con pulpa.
  \item \textbf{Dispersiones coloidales:} Partículas intermedias (1--1\,000\,nm). No se ven a simple vista, pero dispersan la luz (efecto Tyndall). No sedimentan fácilmente. Ej: leche, niebla, gelatina, humo.
  \item \textbf{Dispersiones finas (soluciones):} Partículas a nivel molecular o iónico ($< 1\;\text{nm}$). Completamente homogéneas. No dispersan la luz ni sedimentan. Ej: agua salada, alcohol en agua.
\end{itemize}
\end{cuadrosolucion}

%\newpage

% ============================================================
\subsection{Ejercicios: Clasificación de la Materia}
% ============================================================

\begin{enumerate}[leftmargin=*, resume]
  \item \facil{} Clasifica cada uno como elemento, compuesto, mezcla homogénea o mezcla heterogénea:
  \begin{multicols}{2}
  \begin{enumerate}[label=\alph*)]
    \item Agua destilada
    \item Aire
    \item Oro puro
    \item Ensalada
    \item Acero inoxidable
    \item Dióxido de carbono
  \end{enumerate}
  \end{multicols}

  \item \facil{} ¿Cuál es la diferencia entre un elemento y un compuesto?

  \item \facil{} ¿Cuál es la diferencia entre una mezcla homogénea y una heterogénea?

  \item \facil{} Da dos ejemplos de mezclas que encuentres en tu cocina. Indica si son homogéneas o heterogéneas.

  \item \medio{} Un vaso contiene agua con azúcar completamente disuelta. Otro vaso contiene agua con arena. ¿En cuál tenemos una mezcla homogénea y en cuál una heterogénea? ¿Cómo podrías demostrarlo experimentalmente?

  \item \medio{} ¿El bronce (aleación de cobre y estaño) es una mezcla o un compuesto? Justifica.

  \item \medio{} Clasifica según la clasificación granulométrica (dispersión grosera, coloidal o solución):
  \begin{enumerate}[label=\alph*)]
    \item Café con leche
    \item Agua de mar
    \item Jugo de naranja con pulpa
    \item Niebla
  \end{enumerate}

  \item \medio{} ¿Qué es el efecto Tyndall? ¿Cómo podrías usarlo para distinguir una solución de una dispersión coloidal?

  \item \dificil{} El agua salada se puede separar por evaporación, pero el agua pura (H$_2$O) no se puede separar por evaporación ni por filtración. ¿Significa esto que el agua no se puede descomponer? Si se puede, ¿por qué método?

  \item \dificil{} ¿Por qué los componentes de una mezcla mantienen sus propiedades individuales, pero los de un compuesto no? Usa el agua (H$_2$O) y la mezcla de hidrógeno y oxígeno gaseosos como ejemplo.

  \item \dificil{} Un estudiante afirma: «La leche es una mezcla homogénea porque se ve uniforme.» ¿Está en lo correcto? Argumenta usando la clasificación granulométrica.

  \item \dificil{} Diseña un diagrama de flujo (secuencia lógica de preguntas sí/no) que permita clasificar cualquier muestra de materia como: elemento, compuesto, mezcla homogénea o mezcla heterogénea.
\end{enumerate}

\newpage

% ============================================================
\section{Métodos de Separación de Mezclas}
% ============================================================

\subsection{Fundamento general}

Para separar los componentes de una mezcla aprovechamos las \textbf{diferencias en las propiedades físicas} de cada componente: diferencias de tamaño, densidad, punto de ebullición, propiedades magnéticas, etc.

\begin{cuadroteoria}
\textbf{Principales métodos de separación:}

\begin{enumerate}[leftmargin=*, label=\arabic*.]
  \item \textbf{Filtración:} Separa un sólido de un líquido usando un filtro (papel de filtro, tela). El sólido queda retenido y el líquido pasa. \textit{Propiedad aprovechada:} diferencia de tamaño de partícula. \textit{Ejemplo:} separar arena del agua.

  \item \textbf{Decantación:} Separa líquidos inmiscibles o sólidos de líquidos por diferencia de densidad, dejando reposar la mezcla. \textit{Propiedad:} diferencia de densidad. \textit{Ejemplo:} separar agua y aceite.

  \item \textbf{Evaporación:} Se calienta la mezcla para que el líquido se evapore y el sólido quede. \textit{Propiedad:} diferencia de puntos de ebullición. \textit{Ejemplo:} obtener sal del agua salada.

  \item \textbf{Destilación:} Se calienta la mezcla, se recoge el vapor y se condensa para separar líquidos miscibles. \textit{Propiedad:} diferencia de puntos de ebullición. \textit{Ejemplo:} obtener agua destilada, separar alcohol del agua.

  \item \textbf{Tamizado:} Se pasa la mezcla por un tamiz o colador para separar sólidos de diferentes tamaños. \textit{Propiedad:} diferencia de tamaño de partícula. \textit{Ejemplo:} separar piedras de arena.

  \item \textbf{Imantación:} Se usa un imán para separar materiales ferromagnéticos. \textit{Propiedad:} magnetismo. \textit{Ejemplo:} separar limaduras de hierro de arena.
\end{enumerate}
\end{cuadroteoria}

\begin{figure}[h!]
  \centering
  \begin{tikzpicture}[x=1.3cm, y=1.3cm]
    % Border style
    \tikzset{
      box/.style={draw=mwprimario!40, fill=mwgrisclaro!50, rounded corners, minimum width=4.5cm, minimum height=4cm, align=center}
    }
    
    % 1. FILTRACIÓN (Row 1, Col 1)
    \node[box] at (-4.8, 2.2) {};
    \node[font=\bfseries\small, text=mwprimario] at (-4.8, 3.8) {Filtración};
    \begin{scope}[shift={(-4.8, 2.2)}]
      % Funnel
      \draw[thick, gray] (-0.6, 0.6) -- (0.6, 0.6) -- (0.15, 0) -- (0.15, -0.5) -- (-0.15, -0.5) -- (-0.15, 0) -- cycle;
      % Filter paper
      \draw[thick, draw=mwrojonaranja, fill=mwrojonaranja!10] (-0.55, 0.55) -- (0.55, 0.55) -- (0, 0.05) -- cycle;
      % Beaker
      \draw[thick, gray] (-0.5, -0.6) -- (-0.5, -1.3) -- (0.5, -1.3) -- (0.5, -0.6);
      % Drips
      \fill[mwprimario] (0, -0.2) circle (1.5pt);
      \fill[mwprimario] (0, -0.6) circle (1.5pt);
      \fill[mwprimario] (0, -1.0) circle (1.5pt);
      \fill[mwprimario!12] (-0.45, -1.3) rectangle (0.45, -1.0);
      \draw[thick, gray] (-0.45, -1.0) -- (0.45, -1.0);
      % Label info
      \node[below=1.5cm, font=\tiny, align=center] at (0,0) {Separa sólidos insolubles\\de líquidos (ej. arena y agua).};
    \end{scope}

    % 2. DECANTACIÓN (Row 1, Col 2)
    \node[box] at (0, 2.2) {};
    \node[font=\bfseries\small, text=mwprimario] at (0, 3.8) {Decantación};
    \begin{scope}[shift={(0, 2.2)}]
      % Funnel
      \draw[thick, gray] (-0.6, 0.8) -- (0.6, 0.8) -- (0.15, 0) -- (0.15, -0.5) -- (-0.15, -0.5) -- (-0.15, 0) -- cycle;
      % Liquids
      \fill[mwprimario!12] (-0.48, 0.6) -- (0.48, 0.6) -- (0.15, 0) -- (0.15, -0.3) -- (-0.15, -0.3) -- (-0.15, 0) -- cycle;
      \fill[mwmagenta!12] (-0.55, 0.75) -- (0.55, 0.75) -- (0.48, 0.6) -- (-0.48, 0.6) -- cycle;
      % Separation line
      \draw[thick, draw=mwmagenta] (-0.28, 0.2) -- (0.28, 0.2);
      % Stopcock
      \draw[thick, fill=gray] (-0.2, -0.45) rectangle (0.2, -0.35);
      % Beaker
      \draw[thick, gray] (-0.4, -0.9) -- (-0.4, -1.4) -- (0.4, -1.4) -- (0.4, -0.9);
      \node[below=1.5cm, font=\tiny, align=center] at (0,0) {Separa líquidos inmiscibles\\por densidad (ej. agua y aceite).};
    \end{scope}

    % 3. EVAPORACIÓN (Row 1, Col 3)
    \node[box] at (4.8, 2.2) {};
    \node[font=\bfseries\small, text=mwprimario] at (4.8, 3.8) {Evaporación};
    \begin{scope}[shift={(4.8, 2.2)}]
      % Tripod
      \draw[thick, gray] (-0.6, -1.0) -- (-0.3, -0.2) -- (0.3, -0.2) -- (0.6, -1.0);
      \draw[thick, gray] (0, -0.2) -- (0, -1.0);
      % Capsule
      \draw[thick, fill=mwprimario!12] (-0.5, 0) arc (180:360:0.5cm and 0.2cm) -- cycle;
      % Burner
      \draw[thick, fill=gray!40] (-0.2, -1.0) -- (0.2, -1.0) -- (0.1, -0.6) -- (-0.1, -0.6) -- cycle;
      % Flame
      \fill[mwrojonaranja!80] (-0.1, -0.6) to[out=60, in=-120] (0, -0.2) to[out=-60, in=120] (0.1, -0.6) -- cycle;
      % Vapor waves
      \draw[->, >=stealth, gray, dashed] (-0.25, 0.1) to[out=90, in=-90] (-0.2, 0.5);
      \draw[->, >=stealth, gray, dashed] (0.25, 0.1) to[out=90, in=-90] (0.2, 0.5);
      \node[below=1.5cm, font=\tiny, align=center] at (0,0) {Separa sólidos disueltos\\de líquidos (ej. sal de agua).};
    \end{scope}

    % 4. DESTILACIÓN (Row 2, Col 1)
    \node[box] at (-4.8, -2.2) {};
    \node[font=\bfseries\small, text=mwprimario] at (-4.8, -0.6) {Destilación};
    \begin{scope}[shift={(-4.8, -2.2)}]
      % Flask
      \draw[thick, gray, fill=mwprimario!12] (-0.8, -0.3) circle (0.35cm);
      \draw[thick, gray] (-0.85, 0) -- (-0.85, 0.4) -- (-0.75, 0.4) -- (-0.75, 0);
      % Condenser
      \draw[ultra thick, gray] (-0.8, 0.3) -- (0.2, -0.1);
      \draw[thick, draw=mwprimario] (-0.6, 0.35) -- (0.0, 0.1) -- (0.0, -0.1) -- (-0.6, 0.15) -- cycle;
      % Receiving beaker
      \draw[thick, gray] (0.2, -0.2) -- (0.2, -0.7) -- (0.7, -0.7) -- (0.7, -0.2);
      \fill[mwprimario!12] (0.25, -0.7) rectangle (0.65, -0.5);
      \draw[thick, gray] (0.25, -0.5) -- (0.65, -0.5);
      \node[below=1.5cm, font=\tiny, align=center] at (0,0) {Separa líquidos miscibles por\\puntos de ebullición (ej. alcohol/agua).};
    \end{scope}

    % 5. TAMIZADO (Row 2, Col 2)
    \node[box] at (0, -2.2) {};
    \node[font=\bfseries\small, text=mwprimario] at (0, -0.6) {Tamizado};
    \begin{scope}[shift={(0, -2.2)}]
      % Sieve
      \draw[thick, gray] (-0.8, 0.1) -- (-0.8, -0.1) -- (0.8, -0.1) -- (0.8, 0.1);
      \draw[dashed, ultra thick, gray] (-0.7, -0.1) -- (0.7, -0.1);
      % Large particles (above)
      \filldraw[gray] (-0.4, 0.2) circle (4pt);
      \filldraw[gray] (0.3, 0.15) circle (4pt);
      \filldraw[gray] (-0.1, 0.3) circle (4pt);
      % Small particles (passing)
      \filldraw[gray] (-0.2, -0.3) circle (1.5pt);
      \filldraw[gray] (0.2, -0.5) circle (1.5pt);
      \filldraw[gray] (0.4, -0.3) circle (1.5pt);
      \filldraw[gray] (-0.5, -0.4) circle (1.5pt);
      \node[below=1.5cm, font=\tiny, align=center] at (0,0) {Separa sólidos de distintos\\tamaños (ej. arena y piedras).};
    \end{scope}

    % 6. IMANTACIÓN (Row 2, Col 3)
    \node[box] at (4.8, -2.2) {};
    \node[font=\bfseries\small, text=mwprimario] at (4.8, -0.6) {Imantación};
    \begin{scope}[shift={(4.8, -2.2)}]
      % Horseshoe magnet
      \draw[thick, fill=mwrojonaranja] (-0.3, 0.6) -- (-0.3, 0.2) arc (180:360:0.3cm) -- (0.3, 0.6) -- (0.15, 0.6) -- (0.15, 0.2) arc (360:180:0.15cm) -- (-0.15, 0.6) -- cycle;
      % Blue tips (poles)
      \fill[mwprimario] (0.15, 0.45) rectangle (0.3, 0.6);
      \fill[mwprimario] (-0.3, 0.45) rectangle (-0.15, 0.6);
      % Iron filings attracted
      \fill[gray] (-0.1, -0.2) circle (1.5pt);
      \fill[gray] (0.1, -0.2) circle (1.5pt);
      \fill[gray] (-0.2, -0.3) circle (1.5pt);
      \fill[gray] (0.2, -0.3) circle (1.5pt);
      % Sand pile below
      \filldraw[mwnaranja!12, draw=mwnaranja] (-0.8, -0.9) to[out=30, in=150] (0.8, -0.9) -- cycle;
      \node[below=1.5cm, font=\tiny, align=center] at (0,0) {Separa sustancias magnéticas\\de otras (ej. hierro de arena).};
    \end{scope}
  \end{tikzpicture}
  \caption{Métodos de separación de mezclas.}
  \label{fig:separacion}
\end{figure}

\newpage

% ============================================================
\subsection{Ejercicios: Métodos de Separación}
% ============================================================

\begin{enumerate}[leftmargin=*, resume]
  \item \facil{} ¿Qué método usarías para separar arroz de agua? ¿Qué propiedad aprovechas?

  \item \facil{} ¿Qué método usarías para separar agua de aceite? Justifica.

  \item \facil{} Menciona dos mezclas que se puedan separar por evaporación.

  \item \facil{} ¿Cuál es la diferencia fundamental entre evaporación y destilación?

  \item \medio{} Tienes una mezcla de arena, sal y agua. Diseña un procedimiento paso a paso para separar los tres componentes, indicando qué método usas en cada paso.

  \item \medio{} ¿Qué método usarías para separar alcohol de agua? ¿Por qué no funcionaría la filtración?

  \item \medio{} Tienes una mezcla de limaduras de hierro, arena fina y sal. ¿En qué orden aplicarías los métodos de separación y cuáles serían? Justifica cada paso.

  \item \medio{} Relaciona cada mezcla con su método de separación más adecuado:
  \begin{center}
  \small
  \begin{tabular}{ll}
  \toprule
  \textbf{Mezcla} & \textbf{Método} \\
  \midrule
  a) Piedras y arena & 1) Decantación \\
  b) Agua y aceite & 2) Destilación \\
  c) Agua salada & 3) Tamizado \\
  d) Agua y alcohol & 4) Evaporación \\
  \bottomrule
  \end{tabular}
  \end{center}

  \item \dificil{} ¿Por qué la destilación funciona para separar agua y alcohol pero la evaporación simple no sería adecuada si queremos conservar ambos componentes?

  \item \dificil{} Diseña un diagrama de flujo para separar la siguiente mezcla: limaduras de hierro + arena + sal + agua. Indica en cada paso qué método se usa, qué componente se separa y qué queda.

  \item \dificil{} El petróleo es una mezcla compleja de hidrocarburos. En las refinerías se usa un proceso llamado «destilación fraccionada» para separar sus componentes. Investiga brevemente: ¿en qué se diferencia la destilación fraccionada de la destilación simple?

  \item \dificil{} Un estudiante tiene una mezcla heterogénea de agua, virutas de madera y clips metálicos. Propone el siguiente plan: «Primero filtro, luego uso un imán.» ¿Funcionaría? ¿Qué mejoras le sugerirías y por qué?
\end{enumerate}

\newpage

% ============================================================
\section{Clave de Respuestas}
% ============================================================

\subsection*{Método Científico y Laboratorio (Ejercicios 1--12)}

\begin{enumerate}[leftmargin=*, label=\textbf{\arabic*.}]
  \item Orden: Observación → Planteamiento del problema → Hipótesis → Experimentación → Conclusión.

  \item La observación es un hecho percibido (ej: «las plantas en la ventana crecen más»). La hipótesis es una explicación provisional (ej: «las plantas crecen más porque reciben más luz solar»).

  \item a) Volumétrico. b) De sostén. c) Recipiente. d) De uso específico.

  \item Usar bata y lentes de seguridad; nunca probar ni tocar sustancias químicas directamente.

  \item Hipótesis: «La luz solar favorece el crecimiento de las plantas.» Experimento: plantar semillas idénticas en macetas iguales; unas en la ventana (con luz) y otras en un cuarto oscuro. Variable independiente: cantidad de luz. Variable dependiente: crecimiento de la planta. Variables controladas: tipo de planta, cantidad de agua, tipo de tierra.

  \item a) Pipeta o probeta graduada. b) Mechero de Bunsen y vaso de precipitados (o tubo de ensayo). c) Balanza. d) Embudo con papel de filtro.

  \item Algunos vapores químicos son tóxicos o irritantes. La técnica correcta es el «abanico»: a unos 20\,cm del frasco, agitar suavemente la mano para dirigir una pequeña cantidad de vapor hacia la nariz.

  \item Lavoisier demostró que la masa se conserva en las reacciones químicas (Ley de Conservación de la Masa) y refutó la teoría del flogisto, estableciendo las bases de la química cuantitativa.

  \item Hipótesis: «El azúcar se disuelve más rápido en agua caliente que en agua fría.» Materiales: 4 vasos, agua a 4 temperaturas diferentes, azúcar (misma cantidad), cronómetro, termómetro. Procedimiento: medir $200\;\text{mL}$ de agua a cada temperatura, agregar $10\;\text{g}$ de azúcar simultáneamente, agitar con la misma intensidad y cronometrar hasta disolución completa. Registrar en tabla: temperatura vs. tiempo de disolución.

  \item Sí. Si la hipótesis resulta falsa, el experimento no fracasó: nos dio información valiosa que descarta una posible explicación y nos acerca a la correcta. Toda hipótesis refutada es un avance científico.

  \item No es compatible. En ciencia, ninguna hipótesis se puede «probar definitivamente verdadera»; solo puede ser \emph{no refutada hasta ahora}. La falsabilidad (Karl Popper) establece que una hipótesis científica debe poder ser, en principio, refutada por evidencia empírica.

  \item Medir volumen con probeta graduada, masa con balanza, calcular densidad $d = m/V$, comparar con tablas. Para evaluar si es pura o mezcla: verificar si tiene punto de ebullición constante (sustancia pura) o variable (mezcla) usando termómetro y mechero.
\end{enumerate}

\subsection*{Materia y Propiedades (Ejercicios 13--24)}

\begin{enumerate}[leftmargin=*, label=\textbf{\arabic*.}]
  \setcounter{enumi}{12}
  \item a) E \quad b) I \quad c) E \quad d) I \quad e) I \quad f) E

  \item La masa es la cantidad de materia (constante, escalar, en kg). El peso es la fuerza gravitatoria (varía según $g$, vectorial, en N).

  \item Tierra: $W = 5 \times 9{,}8 = 49\;\text{N}$. Luna: $W = 5 \times 1{,}62 = 8{,}1\;\text{N}$.

  \item Porque no depende de la cantidad: $1\;\text{cm}^3$ de agua tiene $d = 1\;\text{g/cm}^3$ y $1\;\text{L}$ de agua también tiene $d = 1\;\text{g/cm}^3$.

  \item $d_1 = 100/37 \approx 2{,}70\;\text{g/cm}^3$, $d_2 = 250/92{,}6 \approx 2{,}70\;\text{g/cm}^3$. Sí, son del mismo material (probablemente aluminio).

  \item a) F: la masa es constante. b) F: solo a presión de 1\,atm; a mayor altitud hierve a menor temperatura. c) F: el volumen es extensivo, no identifica. d) F: se mide en newtons.

  \item Cambian: masa, volumen, peso (extensivas). Se mantienen: densidad, color, punto de fusión, conductividad (intensivas).

  \item Sí, el aire tiene masa (se puede pesar inflando un globo) y ocupa espacio (llena una habitación). Cumple la definición de materia.

  \item $m = W/g = 147/9{,}8 = 15\;\text{kg}$. En el otro planeta: $W = 15 \times 3{,}5 = 52{,}5\;\text{N}$.

  \item En la vida cotidiana se confunde masa con peso. Deberíamos decir «mi masa es 70\,kg» o «peso 686\,N». La confusión persiste porque en la superficie terrestre la gravedad es casi constante, así que masa y peso son proporcionales.

  \item $d_1 = 50/18{,}5 = 2{,}70$, $d_2 = 120/44{,}4 = 2{,}70$, $d_3 = 200/74{,}1 = 2{,}70\;\text{g/cm}^3$. Todas iguales: mismo material. Color plateado + $d = 2{,}70$ + punto de fusión $660\;°$C = \textbf{aluminio}.

  \item $m = 784/9{,}8 = 80\;\text{kg}$. Marte: $80 \times 3{,}72 = 297{,}6\;\text{N}$. Júpiter: $80 \times 24{,}79 = 1\,983{,}2\;\text{N}$. En Júpiter pesaría más que en la Tierra.
\end{enumerate}

\subsection*{Clasificación de la Materia (Ejercicios 25--36)}

\begin{enumerate}[leftmargin=*, label=\textbf{\arabic*.}]
  \setcounter{enumi}{24}
  \item a) Compuesto. b) Mezcla homogénea. c) Elemento. d) Mezcla heterogénea. e) Mezcla homogénea (aleación). f) Compuesto.

  \item Un elemento no se puede descomponer químicamente. Un compuesto está formado por dos o más elementos en proporción fija y se puede descomponer.

  \item En la homogénea no se distinguen los componentes (una sola fase). En la heterogénea sí se distinguen (dos o más fases).

  \item Ejemplos: café con leche (homogénea), ensalada de frutas (heterogénea).

  \item Agua con azúcar: homogénea (transparente, una fase). Agua con arena: heterogénea (se ven las partículas, sedimentan). Se demuestra dejando reposar: la arena se deposita, el azúcar no.

  \item Es una mezcla (homogénea/aleación). Los componentes no reaccionaron químicamente y se pueden separar por métodos físicos; mantienen sus propiedades individuales a nivel atómico.

  \item a) Coloidal. b) Solución. c) Dispersión grosera. d) Coloidal.

  \item Es la dispersión de un haz de luz al pasar a través de un coloide (las partículas dispersan la luz). En una solución verdadera la luz pasa sin dispersarse. Se puede usar una linterna para distinguirlos.

  \item El agua sí se puede descomponer, pero por \textbf{electrólisis} (método químico, no físico): $2\text{H}_2\text{O} \xrightarrow{\text{electricidad}} 2\text{H}_2 + \text{O}_2$.

  \item En una mezcla, las sustancias solo se juntan sin unirse químicamente. En un compuesto, los átomos se unen por enlaces químicos formando una sustancia nueva. El H$_2$ es inflamable y el O$_2$ alimenta la combustión, pero combinados forman H$_2$O, que apaga el fuego.

  \item No está del todo correcto. La leche es una \textbf{dispersión coloidal} (mezcla heterogénea a nivel microscópico). Aunque parece uniforme, contiene gotas de grasa dispersas que se evidencian con el efecto Tyndall.

  \item ¿Se distinguen componentes a simple vista? → Sí: mezcla heterogénea. No → ¿Se puede separar por métodos físicos? → Sí: mezcla homogénea. No → ¿Se puede descomponer por métodos químicos? → Sí: compuesto. No: elemento.
\end{enumerate}

\subsection*{Métodos de Separación (Ejercicios 37--48)}

\begin{enumerate}[leftmargin=*, label=\textbf{\arabic*.}]
  \setcounter{enumi}{36}
  \item Filtración. Se aprovecha la diferencia de tamaño: el arroz queda en el filtro y el agua pasa.

  \item Decantación (embudo de decantación). Se aprovecha que son inmiscibles y tienen diferente densidad.

  \item Agua salada (la sal queda al evaporar); agua azucarada.

  \item En la evaporación, se pierde el líquido (solo se recupera el sólido). En la destilación, se recuperan \textbf{ambos} componentes: el vapor se condensa y se recoge.

  \item 1) Filtrar: separa la arena (queda en el filtro) del agua salada (pasa). 2) Evaporar: el agua se evapora y la sal queda en la cápsula. Resultado: arena, sal y agua separados (el agua se pierde si solo evaporamos; para recuperarla usaríamos destilación).

  \item Destilación, porque el alcohol ($T_{\text{eb}} = 78\;°$C) y el agua ($T_{\text{eb}} = 100\;°$C) son miscibles. La filtración no sirve porque ambos son líquidos disueltos entre sí.

  \item 1) Imantación: se pasa un imán para extraer las limaduras de hierro. 2) Disolver en agua: la sal se disuelve, la arena no. 3) Filtración: se separa la arena del agua salada. 4) Evaporación: se evapora el agua y queda la sal.

  \item a)-3 (tamizado). b)-1 (decantación). c)-4 (evaporación). d)-2 (destilación).

  \item En la evaporación simple, el vapor del alcohol se pierde en el aire. La destilación condensa el vapor mediante un refrigerante, permitiendo recuperar el alcohol líquido por separado.

  \item Paso 1: Imantación → se extraen limaduras de hierro. Paso 2: Filtración → se separa la arena (sólida) del agua salada (líquida). Paso 3: Evaporación o destilación → se separa la sal del agua.

  \item En la destilación simple hay un solo condensador y se separan dos componentes con puntos de ebullición muy diferentes. En la fraccionada, se usa una columna de fraccionamiento que permite separar múltiples componentes con puntos de ebullición cercanos.

  \item Funcionaría parcialmente. Sugerencias: 1) Primero usar imantación (separar clips); es más eficiente hacerlo antes de filtrar. 2) Luego filtrar para separar virutas de madera del agua (las virutas también flotan, así que la decantación podría ser más sencilla). Plan mejorado: imantación → decantación (virutas flotan) → agua pura.
\end{enumerate}

\end{document}
